\section{Derived tensor product}
\label{section-derived-tensor-product}

\noindent
We can construct the derived tensor product in greater generality.
In fact, it turns out that the boundedness assumptions are not
necessary, provided we choose K-flat resolutions.

\begin{definition}
\label{definition-K-flat}
Let $R$ be a ring. A complex $K^\bullet$ is called {\it K-flat}
if for every acyclic complex $M^\bullet$ the total complex
$\text{Tot}(M^\bullet \otimes_R K^\bullet)$ is acyclic.
\end{definition}

\begin{lemma}
\label{lemma-K-flat-quasi-isomorphism}
Let $R$ be a ring. Let $K^\bullet$ be a K-flat complex.
Then the functor
$$
K(R) \longrightarrow K(R), \quad
L^\bullet \longmapsto \text{Tot}(L^\bullet \otimes_R K^\bullet)
$$
transforms quasi-isomorphisms into quasi-isomorphisms.
\end{lemma}

\begin{proof}
Follows from
Lemma \ref{lemma-derived-tor-exact}
and the fact that quasi-isomorphisms in $K(R)$
are characterized by having acyclic cones.
\end{proof}

\begin{lemma}
\label{lemma-base-change-K-flat}
Let $R \to R'$ be a ring map. If $K^\bullet$ is a K-flat complex
of $R$-modules, then $K^\bullet \otimes_R R'$ is a K-flat complex
of $R'$-modules.
\end{lemma}

\begin{proof}
Follows from the definitions and the fact that
$(K^\bullet \otimes_R R') \otimes_{R'} L^\bullet =
K^\bullet \otimes_R L^\bullet$ for any complex
$L^\bullet$ of $R'$-modules.
\end{proof}

\begin{lemma}
\label{lemma-tensor-product-K-flat}
Let $R$ be a ring. If $K^\bullet$, $L^\bullet$ are K-flat complexes
of $R$-modules, then $\text{Tot}(K^\bullet \otimes_R L^\bullet)$ is a
K-flat complex of $R$-modules.
\end{lemma}

\begin{proof}
Follows from the isomorphism
$$
\text{Tot}(M^\bullet \otimes_R \text{Tot}(K^\bullet \otimes_R L^\bullet))
=
\text{Tot}(\text{Tot}(M^\bullet \otimes_R K^\bullet) \otimes_R L^\bullet)
$$
and the definition.
\end{proof}

\begin{lemma}
\label{lemma-K-flat-two-out-of-three}
Let $R$ be a ring. Let $(K_1^\bullet, K_2^\bullet, K_3^\bullet)$ be
a distinguished triangle in $K(R)$. If two out of three
of $K_i^\bullet$ are K-flat, so is the third.
\end{lemma}

\begin{proof}
Follows from
Lemma \ref{lemma-derived-tor-exact}
and the fact that in a distinguished triangle in
$K(R)$ if two out of three are acyclic, so is the third.
\end{proof}

\begin{lemma}
\label{lemma-K-flat-two-out-of-three-ses}
Let $R$ be a ring. Let
$0 \to K_1^\bullet \to K_2^\bullet \to K_3^\bullet \to 0$ be
a short exact sequence of complexes. If $K_3^n$ is flat for
all $n \in \mathbf{Z}$ and two out of three
of $K_i^\bullet$ are K-flat, so is the third.
\end{lemma}

\begin{proof}
Let $L^\bullet$ be a complex of $R$-modules. Then
$$
0 \to
\text{Tot}(L^\bullet \otimes_R K_1^\bullet) \to
\text{Tot}(L^\bullet \otimes_R K_2^\bullet) \to
\text{Tot}(L^\bullet \otimes_R K_3^\bullet) \to 0
$$
is a short exact sequence of complexes. Namely, for each
$n, m$ the sequence of modules
$0 \to L^n \otimes_R K_1^m \to
L^n \otimes_R K_2^m \to
L^n \otimes_R K_3^m \to 0$
is exact by Algebra, Lemma \ref{algebra-lemma-flat-tor-zero}
and the sequence of complexes is a direct sum of these.
Thus the lemma follows from this
and the fact that in a short exact sequence of complexes
if two out of three are acyclic, so is the third.
\end{proof}

\begin{lemma}
\label{lemma-derived-tor-quasi-isomorphism}
Let $R$ be a ring. Let $P^\bullet$ be a bounded above complex of
flat $R$-modules. Then $P^\bullet$ is K-flat.
\end{lemma}

\begin{proof}
Let $L^\bullet$ be an acyclic complex of $R$-modules.
Let $\xi \in H^n(\text{Tot}(L^\bullet \otimes_R P^\bullet))$.
We have to show that $\xi = 0$.
Since $\text{Tot}^n(L^\bullet \otimes_R P^\bullet)$ is a direct
sum with terms $L^a \otimes_R P^b$ we see that $\xi$ comes from
an element in $H^n(\text{Tot}(\tau_{\leq m}L^\bullet \otimes_R P^\bullet))$
for some $m \in \mathbf{Z}$. Since $\tau_{\leq m}L^\bullet$ is also
acyclic we may replace $L^\bullet$ by $\tau_{\leq m}L^\bullet$.
Hence we may assume that $L^\bullet$ is bounded above.
In this case the spectral sequence of
Homology, Lemma \ref{homology-lemma-first-quadrant-ss}
has
$$
{}'E_1^{p, q} = H^p(L^\bullet \otimes_R P^q)
$$
which is zero as $P^q$ is flat and $L^\bullet$ acyclic. Hence
$H^*(\text{Tot}(L^\bullet \otimes_R P^\bullet)) = 0$.
\end{proof}

\noindent
In the following lemma by a colimit of a system of complexes we mean
the termwise colimit.

\begin{lemma}
\label{lemma-colimit-K-flat}
Let $R$ be a ring.
Let $K_1^\bullet \to K_2^\bullet \to \ldots$
be a system of K-flat complexes.
Then $\colim_i K_i^\bullet$ is K-flat.
More generally any filtered colimit of K-flat complexes
is K-flat.
\end{lemma}

\begin{proof}
Because we are taking termwise colimits we have
$$
\colim_i \text{Tot}(M^\bullet \otimes_R K_i^\bullet) =
\text{Tot}(M^\bullet \otimes_R \colim_i K_i^\bullet)
$$
by Algebra, Lemma \ref{algebra-lemma-tensor-products-commute-with-limits}.
Hence the lemma follows from the fact that filtered colimits are
exact, see Algebra, Lemma \ref{algebra-lemma-directed-colimit-exact}.
\end{proof}

\begin{lemma}
\label{lemma-universally-acyclic-K-flat}
Let $R$ be a ring. Let $K^\bullet$ be a complex of $R$-modules.
If $K^\bullet \otimes_R M$ is acyclic for all finitely presented
$R$-modules $M$, then $K^\bullet$ is K-flat.
\end{lemma}

\begin{proof}
We will use repeatedly that tensor product commute with colimits
(Algebra, Lemma \ref{algebra-lemma-tensor-products-commute-with-limits}).
Thus we see that $K^\bullet \otimes_R M$ is acyclic for
any $R$-module $M$, because any $R$-module is a filtered colimit
of finitely presented $R$-modules $M$, see
Algebra, Lemma \ref{algebra-lemma-module-colimit-fp}.
Let $M^\bullet$ be an acyclic complex of $R$-modules.
We have to show that $\text{Tot}(M^\bullet \otimes_R K^\bullet)$
is acyclic. Since $M^\bullet = \colim \tau_{\leq n} M^\bullet$ (termwise
colimit) we have
$$
\text{Tot}(M^\bullet \otimes_R K^\bullet) =
\colim \text{Tot}(\tau_{\leq n} M^\bullet \otimes_R K^\bullet)
$$
with truncations as in Homology, Section \ref{homology-section-truncations}.
As filtered colimits are exact
(Algebra, Lemma \ref{algebra-lemma-directed-colimit-exact})
we may replace $M^\bullet$ by $\tau_{\leq n}M^\bullet$ and
assume that $M^\bullet$ is bounded above.
In the bounded above case, we can write
$M^\bullet = \colim \sigma_{\geq -n} M^\bullet$
where the complexes $\sigma_{\geq -n} M^\bullet$ are bounded
but possibly no longer acyclic.
Arguing as above we reduce to the case where $M^\bullet$
is a bounded complex.
Finally, for a bounded complex $M^a \to \ldots \to M^b$
we can argue by induction on the length $b - a$ of the complex.
The case $b - a = 1$ we have seen above.
For $b - a > 1$ we consider the split short exact sequence
of complexes
$$
0 \to \sigma_{\geq a + 1}M^\bullet \to M^\bullet \to M^a[-a] \to 0
$$
and we apply Lemma \ref{lemma-derived-tor-exact}
to do the induction step. Some details omitted.
\end{proof}

\begin{lemma}
\label{lemma-K-flat-resolution}
Let $R$ be a ring. For any complex $M^\bullet$ there exists a
K-flat complex $K^\bullet$ whose terms are flat $R$-modules
and a quasi-isomorphism $K^\bullet \to M^\bullet$ which is termwise
surjective.
\end{lemma}

\begin{proof}
Let $\mathcal{P} \subset \Ob(\text{Mod}_R)$ be the
class of flat $R$-modules. By
Derived Categories, Lemma \ref{derived-lemma-special-direct-system}
there exists a system
$K_1^\bullet \to K_2^\bullet \to \ldots$
and a diagram
$$
\xymatrix{
K_1^\bullet \ar[d] \ar[r] &
K_2^\bullet \ar[d] \ar[r] & \ldots \\
\tau_{\leq 1}M^\bullet \ar[r] &
\tau_{\leq 2}M^\bullet \ar[r] & \ldots
}
$$
with the properties (1), (2), (3) listed in that lemma.
These properties imply each complex $K_i^\bullet$ is a bounded
above complex of flat modules. Hence $K_i^\bullet$ is K-flat by
Lemma \ref{lemma-derived-tor-quasi-isomorphism}.
The induced map $\colim_i K_i^\bullet \to M^\bullet$
is a quasi-isomorphism and termwise surjective by construction. The complex
$\colim_i K_i^\bullet$ is K-flat by
Lemma \ref{lemma-colimit-K-flat}.
The terms $\colim K_i^n$ are flat because filtered colimits of
flat modules are flat, see
Algebra, Lemma \ref{algebra-lemma-colimit-flat}.
\end{proof}

\begin{remark}
\label{remark-P-resolution}
In fact, we can do better than Lemma \ref{lemma-K-flat-resolution}.
Namely, we can find a quasi-isomorphism
$P^\bullet \to M^\bullet$ where $P^\bullet$ is a complex of $R$-modules
endowed with a filtration
$$
0 = F_{-1}P^\bullet \subset F_0P^\bullet \subset
F_1P^\bullet \subset \ldots \subset P^\bullet
$$
by subcomplexes such that
\begin{enumerate}
\item $P^\bullet = \bigcup F_pP^\bullet$,
\item the inclusions $F_iP^\bullet \to F_{i + 1}P^\bullet$
are termwise split injections,
\item the quotients $F_{i + 1}P^\bullet/F_iP^\bullet$ are isomorphic to direct
sums of shifts $R[k]$ (as complexes, so differentials are zero).
\end{enumerate}
This will be shown in
Differential Graded Algebra, Lemma \ref{dga-lemma-resolve}
(or you can argue as in the proof of
Lemma \ref{lemma-resolution-filtered-complex}).
Moreover, given such a complex we obtain a distinguished triangle
$$
\bigoplus F_iP^\bullet \to \bigoplus F_iP^\bullet \to M^\bullet
\to \bigoplus F_iP^\bullet[1]
$$
in $D(R)$. Using this we can sometimes reduce statements about general
complexes to statements about $R[k]$ (this of course only works if the
statement is preserved under taking direct sums). More precisely, let
$T$ be a property of objects of $D(R)$. Suppose that
\begin{enumerate}
\item if $K_i \in D(R)$, $i \in I$ is a family of objects with
$T(K_i)$ for all $i \in I$, then $T(\bigoplus K_i)$,
\item if $K \to L \to M \to K[1]$ is a distinguished triangle and
$T$ holds for two, then $T$ holds for the third object,
\item $T(R[k])$ holds for all $k$.
\end{enumerate}
Then $T$ holds for all objects of $D(R)$.
\end{remark}

\begin{lemma}
\label{lemma-derived-tor-quasi-isomorphism-other-side}
Let $R$ be a ring. Let
$\alpha : P^\bullet \to Q^\bullet$ be a quasi-isomorphism of
K-flat complexes of $R$-modules. For every complex $L^\bullet$
of $R$-modules the induced map
$$
\text{Tot}(\text{id}_L \otimes \alpha) :
\text{Tot}(L^\bullet \otimes_R P^\bullet)
\longrightarrow
\text{Tot}(L^\bullet \otimes_R Q^\bullet)
$$
is a quasi-isomorphism.
\end{lemma}

\begin{proof}
Choose a quasi-isomorphism $K^\bullet \to L^\bullet$ with
$K^\bullet$ a K-flat complex, see
Lemma \ref{lemma-K-flat-resolution}.
Consider the commutative diagram
$$
\xymatrix{
\text{Tot}(K^\bullet \otimes_R P^\bullet) \ar[r] \ar[d] &
\text{Tot}(K^\bullet \otimes_R Q^\bullet) \ar[d] \\
\text{Tot}(L^\bullet \otimes_R P^\bullet) \ar[r] &
\text{Tot}(L^\bullet \otimes_R Q^\bullet)
}
$$
The result follows as by
Lemma \ref{lemma-K-flat-quasi-isomorphism}
the vertical arrows and the top horizontal arrow are quasi-isomorphisms.
\end{proof}

\noindent
Let $R$ be a ring.
Let $M^\bullet$ be an object of $D(R)$.
Choose a K-flat resolution $K^\bullet \to M^\bullet$, see
Lemma \ref{lemma-K-flat-resolution}.
By
Lemmas \ref{lemma-derived-tor-homotopy} and \ref{lemma-derived-tor-exact}
we obtain an exact functor of triangulated categories
$$
K(R) \longrightarrow K(R), \quad
L^\bullet \longmapsto \text{Tot}(L^\bullet \otimes_R K^\bullet)
$$
By
Lemma \ref{lemma-K-flat-quasi-isomorphism}
this functor induces a functor $D(R) \to D(R)$ simply because
$D(R)$ is the localization of $K(R)$ at quasi-isomorphism.
As the category of $K$-flat resolutions of $M^\bullet$ is
cofiltered and as we have
Lemma \ref{lemma-derived-tor-quasi-isomorphism-other-side}
the resulting functor (up to isomorphism)
does not depend on the choice of $K^\bullet$.

\begin{definition}
\label{definition-derived-tor}
Let $R$ be a ring. Let $M^\bullet$ be an object of $D(R)$.
The {\it derived tensor product}
$$
- \otimes_R^{\mathbf{L}} M^\bullet : D(R) \longrightarrow D(R)
$$
is the exact functor of triangulated categories described above.
\end{definition}

\noindent
This functor extends the functor (\ref{equation-derived-tensor-module}).
It is clear from our explicit constructions that
there is an isomorphism (involving a choice of signs, see below)
$$
M^\bullet \otimes_R^{\mathbf{L}} L^\bullet
\cong
L^\bullet \otimes_R^{\mathbf{L}} M^\bullet
$$
whenever both $L^\bullet$ and $M^\bullet$ are in $D(R)$.
Hence when we write $M^\bullet \otimes_R^{\mathbf{L}} L^\bullet$
we will usually be agnostic about which variable we are using to
define the derived tensor product with.

\begin{lemma}
\label{lemma-flip-douoble-tensor-product}
Let $R$ be a ring. Let $K^\bullet, L^\bullet$ be complexes of $R$-modules.
There is a canonical isomorphism
$$
K^\bullet \otimes_R^\mathbf{L} L^\bullet \longrightarrow
L^\bullet \otimes_R^\mathbf{L} K^\bullet
$$
functorial in both complexes which uses a sign of $(-1)^{pq}$
for the map $K^p \otimes_R L^q \to L^q \otimes_R K^p$ (see proof
for explanation).
\end{lemma}

\begin{proof}
We may and do replace the complexes by K-flat complexes $K^\bullet$ and
$L^\bullet$ and then we use the commutativity constraint discussed
in Section \ref{section-symmetric-monoidal}.
\end{proof}

\begin{lemma}
\label{lemma-triple-tensor-product}
Let $R$ be a ring. Let $K^\bullet, L^\bullet, M^\bullet$
be complexes of $R$-modules. There is a canonical isomorphism
$$
(K^\bullet \otimes_R^\mathbf{L} L^\bullet) \otimes_R^\mathbf{L} M^\bullet
=
K^\bullet \otimes_R^\mathbf{L} (L^\bullet \otimes_R^\mathbf{L} M^\bullet)
$$
functorial in all three complexes.
\end{lemma}

\begin{proof}
Replace the complexes by K-flat complexes and use the
associativity constraint in Section \ref{section-symmetric-monoidal}.
\end{proof}

\begin{lemma}
\label{lemma-factor-through-K-flat}
Let $R$ be a ring. Let $a : K^\bullet \to L^\bullet$ be a map of complexes
of $R$-modules. If $K^\bullet$ is K-flat, then there exist a complex
$N^\bullet$ and maps of complexes $b : K^\bullet \to N^\bullet$
and $c : N^\bullet \to L^\bullet$ such that
\begin{enumerate}
\item $N^\bullet$ is K-flat,
\item $c$ is a quasi-isomorphism,
\item $a$ is homotopic to $c \circ b$.
\end{enumerate}
If the terms of $K^\bullet$ are flat, then we may choose
$N^\bullet$, $b$, and $c$
such that the same is true for $N^\bullet$.
\end{lemma}

\begin{proof}
We will use that the homotopy category $K(R)$
is a triangulated category, see Derived Categories, Proposition
\ref{derived-proposition-homotopy-category-triangulated}.
Choose a distinguished triangle
$K^\bullet \to L^\bullet \to C^\bullet \to K^\bullet[1]$.
Choose a quasi-isomorphism $M^\bullet \to C^\bullet$ with
$M^\bullet$ K-flat with flat terms, see Lemma \ref{lemma-K-flat-resolution}.
By the axioms of triangulated categories,
we may fit the composition $M^\bullet \to C^\bullet \to K^\bullet[1]$
into a distinguished triangle
$K^\bullet \to N^\bullet \to M^\bullet \to K^\bullet[1]$.
By Lemma \ref{lemma-K-flat-two-out-of-three} we see that $N^\bullet$ is K-flat.
Again using the axioms of triangulated categories,
we can choose a map $N^\bullet \to L^\bullet$ fitting into
the following morphism of distinguished triangles
$$
\xymatrix{
K^\bullet \ar[r] \ar[d] &
N^\bullet \ar[r] \ar[d] &
M^\bullet \ar[r] \ar[d] &
K^\bullet[1] \ar[d] \\
K^\bullet \ar[r] &
L^\bullet \ar[r] &
C^\bullet \ar[r] &
K^\bullet[1]
}
$$
Since two out of three of the arrows are quasi-isomorphisms, so is
the third arrow $N^\bullet \to L^\bullet$ by the long exact sequences
of cohomology associated to these distinguished triangles
(or you can look at the image of this diagram in $D(R)$ and use
Derived Categories, Lemma \ref{derived-lemma-third-isomorphism-triangle}
if you like). This finishes the proof of (1), (2), and (3).
To prove the final assertion, we may choose $N^\bullet$
such that $N^n \cong M^n \oplus K^n$, see
Derived Categories, Lemma
\ref{derived-lemma-improve-distinguished-triangle-homotopy}.
Hence we get the desired flatness
if the terms of $K^\bullet$ are flat.
\end{proof}








